\section{Introduction}

Wind energy has become one of the main sources of low-carbon electricity, and the placement of the turbines, fixed at the planning stage, has a lasting effect on the energy a farm delivers. The wind farm layout optimization problem (WFLOP) addresses this decision. An upstream turbine leaves behind a wake of slower and more turbulent air that reduces the power of the turbines downstream; the wake recovers with distance, but the site is limited and the turbines must keep a minimum distance from each other. The objective is nonlinear, the wakes couple the positions of all turbines, and the feasible region is bounded by many nonconvex spacing constraints, which has made WFLOP a popular test bed for metaheuristics.

Early studies placed turbines on a grid with genetic algorithms (GAs)~\cite{Mosetti1994,Grady2005}. Kusiak and Song~\cite{Kusiak2010} introduced the continuous, wind-distribution-based model that we also use, and it has since been solved with ant colony optimization and particle filtering~\cite{Eroglu2012,Eroglu2013} and with biogeography-based optimization~\cite{Bansal2017}. GA, particle swarm optimization (PSO), differential evolution (DE) and simulated annealing were compared in~\cite{Elkinton2008}, PSO has been applied to reduce wake losses~\cite{Asaah2021}, and the collector and electrical systems of offshore farms have been optimized and reviewed in~\cite{Srikakulapu2018,Hou2019}. Later work used informed GAs~\cite{Ju2019,Bai2022}, deterministic refinement~\cite{Nagpal2021}, variable neighborhood search (VNS) for large offshore farms~\cite{Cazzaro2022}, pseudo-gradients~\cite{Quaeghebeur2021}, Gaussian wake models~\cite{Tao2020} and surrogates~\cite{Yang2023,Wang2024,Li2025}. The IEA Wind Task~37 case study~\cite{Baker2019} showed that, even with a common wake model, different optimizers reach clearly different layouts.

Many of these studies propose a new or modified metaheuristic and compare it with reference methods, among which PSO is one of the most common. Such comparisons are only as informative as the configuration of the references. The particles of PSO converge only if its inertia weight $w$ and acceleration coefficients $c_1$, $c_2$ lie in a known region~\cite{Clerc2002,Poli2009}, and the setting $w=0.7$, $c_1=c_2=2$, which is common in PSO baselines and was used in the earlier version of this study, lies outside it. With the constriction coefficients of Clerc and Kennedy~\cite{Clerc2002}, plain PSO outperforms the salp swarm algorithm (SSA)~\cite{Mirjalili2017}, the Laplacian SSA (LX-SSA)~\cite{Solanki2023} proposed earlier by two of the authors, and an SSA-VNS hybrid that we had developed first (Section~\ref{sec:results}).
% CHECK-FINAL [C12]: main.friedman.avg_rank: PSOC below SSA, LXSSA, SSABV
A baseline outside its convergence region can therefore make a new method look stronger than it is; we report this as a finding about the configuration of baselines, not about any particular study.

Swarms and trajectory methods complement each other. Once a swarm has contracted around its best layout, the remaining evaluations bring small gains, because the final refinement requires small moves of individual turbines along the active spacing constraints, whereas VNS~\cite{Mladenovic1997,Hansen2001} refines one layout efficiently but depends on where it starts. A swarm-then-VNS design helps only if the swarm provides a better start than cheaper alternatives, and SSA did not: SSA-VNS was statistically tied with a control that spends the same evaluations on random layouts (Section~\ref{sec:ablation}).
% CHECK-FINAL [C15]: ablation.contrasts.SSABV-RSVNS: T >= 60, W and L <= 5
A convergent PSO, which contracts around a good region within its share of the budget, is the natural alternative.

We therefore propose PSO-VNS: PSO with constriction coefficients explores the layout space for the first half of the budget, and the basic VNS refines the global best layout of the swarm. Both components are standard, and the gain over a well-configured PSO is modest; the contribution lies in their combination for the constrained WFLOP, in a controlled evaluation of which ingredients matter, and in the finding on the configuration of PSO baselines. We keep the Jensen wake model, Weibull wind data and the Kusiak--Song power model for comparability with the literature and examine these choices separately. Our contributions are as follows:
\begin{enumerate}
\item We propose PSO-VNS, a two-phase algorithm for the continuous WFLOP with explicit boundary and minimum-spacing constraints, in which a convergent PSO and the basic VNS share one penalized objective (Section~\ref{sec:hybrid}).
\item We compare PSO-VNS with PSO, SSA-VNS, SSA, LX-SSA, DE, VNS and a gradient-based multistart SLSQP solver on 68 benchmark cases at equal numbers of objective evaluations with 30 seed-paired runs, and isolate in an ablation the VNS phase, the value of the swarm start (against VNS alone and a random-multistart control), the choice of the first-phase swarm and the Laplace step of LX-SSA (Section~\ref{sec:results}).
\item We show that the PSO setting $w=0.7$, $c_1=c_2=2$ violates the convergence condition of PSO and quantify how much it understates PSO (Sections~\ref{sec:pso} and~\ref{sec:psosetting}).
\item We test the methods with feasibility-preserving initialization, budgets of up to 120,030 evaluations, a 16-turbine block of the Horns Rev~1 offshore farm, the 16- and 36-turbine scenarios of IEA Wind Task~37 Case Study~1 (IEA37), and alternative power-curve, wake, spacing and boundary assumptions (Sections~\ref{sec:beyond} and~\ref{sec:robustness}).
\end{enumerate}

Section~\ref{sec:model} describes the model, Section~\ref{sec:prelim} the swarm methods, and Section~\ref{sec:hybrid} PSO-VNS. Sections~\ref{sec:setup}--\ref{sec:robustness} give the setup and results, Section~\ref{sec:limits} the limitations, and Section~\ref{sec:conclusion} concludes the paper.

